% Part: methods
% Chapter: proofs
% Section: resources

\documentclass[../../../include/open-logic-section]{subfiles}

\begin{document}

\olfileid{mth}{prf}{res}

\section{نورې سرچينې}

په رياضياتو کښې د ثبوتونو د کولو په اړه ډېر کتابونه
شته چې ګټور کېدے شي۔ په ځانګړي ډول \emph{ثبوتونه
څنګه ولولو او وکړو: د رياضيکي فکر د بهيرونو پېژندنه}
\citep{Solow2013} او \emph{څنګه ئې ثابت کړو: يوه منظمه
لار} \citep{Velleman2019} وګورئ۔
\href{http://www.people.vcu.edu/~rhammack/BookOfProof/BookOfProof.pdf}{\emph{د ثبوت کتاب}}
\citep{Hammack2013} او
\href{https://scholarworks.gvsu.edu/books/7/}{\emph{رياضيکي استدلال}}
\citep{Sandstrum2019} د ثبوت په اړه داسې کتابونه دي
چې په انټرنېټ کښې وړيا شته۔ د فلسفې لوستونکي ښايي
\emph{لا په دقت: هغه رياضيات چې د فلسفې د کولو
دپاره ئې اړتيا لرئ} \citep{Steinhart2018} د رياضيکي
استدلال ښه لومړنے لارښود وګڼي۔

په انټرنېټ کښې د ثبوتونو بېلابېل لنډ لارښودونه هم
شته؛ د بېلګې په توګه
\href{https://math.berkeley.edu/~hutching/teach/proofs.pdf}{«د رياضيکي استدلالونو پېژندنه»}
\citep{Hutchings2003} او
\href{https://eugeniacheng.com/wp-content/uploads/2017/02/cheng-proofguide.pdf}{«ثبوتونه څنګه وليکو»}
\citep{Cheng2004}۔

\subsection{هڅوونکې ويډيوګانې}

د کورنۍ دندې د کولو زړه مو نۀ کېږي؟ خپه يئ؟
ښايي دا ويډيوګانې مرسته وکړي!

\begin{itemize}
\item \url{https://www.youtube.com/watch?v=ZXsQAXx_ao0}
\item \url{https://www.youtube.com/watch?v=BQ4yd2W50No}
\item \url{https://www.youtube.com/watch?v=StTqXEQ2l-Y}
\end{itemize}

\end{document}
